\begin{abstract}
Internet of Things sites that cannot depend on a remote service run large
language models (LLMs) as pipelines over heterogeneous edge devices.  A stage failure then loses the request's key-value (KV)
cache.  Existing recovery paths recompute the lost state at prefix-length
cost, replicate every protected stage's KV cache, or reactively reconfigure
and reload the pipeline.  We present KV-CARE (KV Cache Availability
and Recovery at the Edge).  Cross-stage parity folds the protected stages'
KV columns into one parity column and reconstructs a failed stage's KV cache bit-exactly without recomputing the
prefix.  Recovery-aware placement charges backup-weight reservations to each
device's memory while placing layers, so the failed stage's weights have a
device to load on.  On a six-stage Jetson pipeline serving OPT-6.7B,
KV-CARE's recovery latency is flat over the measured failure positions
($5\pm3$\,ms per generated token against $447\pm6$\,ms for recomputation)
and completes in two to four decode steps over that range.  It recovers
about one decode step later than replication while retaining 3.7$\times$
less coordinator KV state (2.7$\times$ over the stages it recovers), and it
restores two simultaneous failures with recovery latency comparable to a
single failure when their backups reside on distinct nodes.  In offline
placement analysis with backups confined to pipeline devices, joint
placement remains feasible across the memory range in which the model fits,
where placement followed by backup assignment does not.  Single-parity
protection changes throughput by $+0.6$\%.  On site, a protected-stage
failure whose recovery contract holds becomes a delay of a few decode steps.
\end{abstract}
